\begin{abstract}
An observer who sees a Hamiltonian ensemble through a Gaussian window of width $\sigma$ finds that the coarse-grained entropy grows only if the window widens faster than a critical rate $R$ of the state. For Gaussian states $R$ stays below the largest strain rate $\kappa$ of the flow by a margin proportional to $\sigma$, the resolution gap. We ask whether every state has a gap of order $\sigma$, at the unstable top of the pendulum ($\kappa = 1$). Numerically, it does not. Ripple strips, states that carry a fine ripple along the unstable direction, reach $(\kappa - R)/\sigma = 1/(2k_0)$ as $\sigma \to 0$ for every ripple wavenumber $k_0$ (in units of $1/\sigma$), so that the smallest gap at a given $\sigma$ tends to zero; between $\sigma = 10^{-24}$ and $10^{-96}$ it follows $1.02/\sqrt{\ln(1/\sigma) - 23.4}$, and every value is confirmed by evaluations that share no code. In the other direction we prove that the gap is controlled by one number, the mean squared wavenumber $\bar x_w$ of the blurred density along the unstable direction, and that the blur limits it logarithmically: under one checkable ratio condition, $\kappa - R \geq c\,\sigma/\sqrt{\ln(1/\sigma) + C}$ with explicit constants. We then remove the ratio condition for product states in any orientation and for every state whose posterior never spreads beyond a fixed multiple of the blur. The proof follows the stable-direction Fisher information as the blur grows and shows that the only way to lose it faster is through valleys of the blurred density, which are paid for in Fisher information. The strips, which realise the smallest gaps we found, numerically form no deep valleys, so the states that come closest to the edge are covered without any ratio condition. For general states with deep valleys the unconditional law is stated as a conjecture. Fine structure can narrow the gap, but only logarithmically. Gaussian mixtures with up to 32 components never go below $0.884\,\sigma$, and the strips beat them only for $\sigma \lesssim 7\times 10^{-6}$.
\end{abstract}
